\فصل{کارهای پیشین}

این فصل پژوهش‌های مرتبط را به‌صورت مقایسه‌ای مرور می‌کند. مفاهیم بازشناسی گفتار، قطعه‌بندی، شبه‌برچسب، تخریب مخابراتی، بهسازی و همجوشی در فصل پیش تعریف شده‌اند؛ بنابراین، در این فصل بررسی می‌شود که مطالعات موجود چه شواهدی ارائه کرده‌اند، در چه دامنه‌هایی ارزیابی شده‌اند و چه پرسش‌هایی را دربارهٔ گفتار فارسی بی‌پاسخ گذاشته‌اند. این مرور در چهار محور سامان یافته است: بازشناسی گفتار فارسی و انتقال چندزبانه، ساخت پیکره از گفتار وب و نظارت ضعیف، بازشناسی مقاوم در کانال مخابراتی، و پیوند بهسازی با همجوشی چندنما. پس از این چهار محور، مطالعات منتخب در جدولی مقایسه و شکاف‌های پژوهشی جمع‌بندی می‌شوند.

\قسمت{بازشناسی گفتار فارسی و انتقال چندزبانه}

\زیرقسمت{از پیکره‌های پایه تا مدل‌های عصبی}

پژوهش دربارهٔ بازشناسی گفتار فارسی از آغاز به اندازه و دامنهٔ پیکره‌های در دسترس وابسته بوده است. فارس‌دَت منبعی پایه برای مطالعهٔ گفتار خوانده‌شده و واحدهای آوایی فارسی فراهم کرد \مرجع{bijankhan1994farsdat}. پس از آن، سامانهٔ نویسا مدل مخلوط گاوسی--مارکوف پنهان، واژه‌نامهٔ تلفظ و مدل زبانی را برای بازشناسی پیوسته با واژگان گسترده یکپارچه کرد \مرجع{sameti2008nevisa,sameti2011large}. اهمیت این آثار تنها به خط مبناهای عددی آن‌ها محدود نمی‌شود؛ این پژوهش‌ها نشان می‌دهند که طراحی واژه‌نامه، پیکره و قرارداد رونویسی فارسی، بخش جدایی‌ناپذیری از طراحی بازشناس است. بااین‌حال، غلبهٔ گفتار خوانده‌شده و شرایط ضبط کنترل‌شده در این منابع، تعمیم مستقیم یافته‌ها به گفتار وب یا تلفنی را محدود می‌کند.

با گسترش مدل‌های عصبی، تمرکز از مهندسی جداگانهٔ اجزای کلاسیک به یادگیری بازنمایی‌های آکوستیکی انتقال یافت. ویسی و حاجی‌مانی چند معماری عمیق را برای فارسی بررسی کردند و با وجود بهبود مدل‌سازی، کمبود پیکره‌های بزرگ و استاندارد را همچنان مانع اصلی دانستند \مرجع{veisi2020persian}. دهقانی و سیدصالحی با استفاده از شبکهٔ پیچشی مکس‌اوت بر محلی‌سازی زمان--فرکانس تمرکز کردند. اباسکوهی و همکاران نیز مسئله‌ای تخصصی‌تر، یعنی ارزیابی گفتار کودکان پیش‌دبستانی، را مطالعه کردند \مرجع{dehghani2023time,abaskohi2022automatic}. این دو مسیر دامنهٔ پژوهش دربارهٔ فارسی را گسترش می‌دهند، اما درعین‌حال نشان می‌دهند که «عملکرد در زبان فارسی» کمیتی یگانه نیست؛ سن، سبک گفتار، گسترهٔ واژگان و محیط ضبط می‌توانند ماهیت مسئله را تغییر دهند.

پیکره‌های چندزبانهٔ جدید، دسترسی به داده و امکان مقایسه را بهبود داده‌اند. کامن‌وویس (\lr{Common Voice}) بر مشارکت داوطلبانه و ضبط جمله‌های ارائه‌شده تکیه دارد، درحالی‌که فلورز (\lr{FLEURS}) مجموعه‌ای موازی و چندزبانه برای ارزیابی بازنمایی‌های گفتاری فراهم می‌کند \مرجع{commonvoice2024,conneau2022fleurs}. این دو منبع برای گسترش پوشش زبانی و ارزیابی میان‌زبانی ارزشمندند، اما جای پیکره‌ای ناهمگون از برنامه‌های طولانی یا تماس‌های واقعی را نمی‌گیرند. تفاوت منابع نیز تنها آکوستیکی نیست: جملهٔ خوانده‌شده معمولاً مرزی روشن، متنی از پیش مشخص و یک گوینده دارد؛ حال آنکه محتوای طولانی ممکن است شامل مکث، اصلاح گفتار، موسیقی و تغییر گوینده باشد. بنابراین، افزایش ساعت‌های داده بدون توجه به نوع گفتار، لزوماً شکاف میان دامنه‌ها را کاهش نمی‌دهد.

\زیرقسمت{پیش‌آموزش و ریزتنظیم چندزبانه}

ویو۲وک ۲ (\lr{wav2vec 2.0}) و ویسپر (\lr{Whisper}) دو رویکرد متفاوت برای کاهش وابستگی به رونویسی در دامنهٔ هدف ارائه کرده‌اند. \lr{wav2vec 2.0} ابتدا ساختار گفتار را با هدفی خودنظارتی می‌آموزد و سپس با داده‌های برچسب‌خورده سازگار می‌شود \مرجع{baevski2020wav2vec}. در مقابل، \lr{Whisper} از نظارت ضعیف در مقیاس بزرگ برای یادگیری هم‌زمان بازشناسی، ترجمه و تشخیص زبان استفاده می‌کند \مرجع{radford2023robust}. رویکرد نخست، جدایی مرحلهٔ یادگیری بازنمایی از نظارت زبانی را برجسته می‌کند؛ رویکرد دوم نشان می‌دهد که داده‌های جفت‌شده اما ناهمگون نیز، اگر در مقیاس بزرگ فراهم شوند، می‌توانند پایه‌ای انتقال‌پذیر بسازند. بااین‌حال، هیچ‌یک به‌تنهایی تضمین نمی‌کنند که دامنهٔ هدف، به‌ویژه کانال تلفنی در یک زبان کم‌منبع، به‌اندازهٔ کافی در داده‌های پیش‌آموزش نمایندگی شده باشد.

لیو و همکاران چند راهبرد ریزتنظیم \lr{Whisper} را در زبان‌های کم‌منبع مقایسه کردند و نشان دادند که انتخاب راهبرد به مقدار داده و دامنهٔ آن وابسته است \مرجع{liu2024whisper}. رحیمی‌پور و همکاران نیز در مطالعه‌ای دربارهٔ فارسی، ریزتنظیم \lr{Whisper} را با استفادهٔ مستقیم از مدل آماده و خط مبناهای دیگر مقایسه کردند و شواهد مثبتی از سازگارسازی مدل گزارش دادند \مرجع{rahimipour2024persian}. وجه مشترک این مطالعات آن است که پیش‌آموزش نقطهٔ شروع بهتری فراهم می‌کند، اما نیاز به داده‌های دامنهٔ هدف را از میان نمی‌برد. محدودیت مشترک آن‌ها برای بازشناسی مقاوم نیز این است که ارزیابی در زبان‌های کم‌منبع معمولاً بر زبان و پیکره متمرکز می‌شود و اثر مستقل کدک، افت بسته یا سبک تماس را آشکار نمی‌کند.

کانفورمر (\lr{Conformer}) با ترکیب مدل‌سازی محلی و سراسری، مسیر مهم دیگری را در پژوهش بازشناسی خودکار گفتار (\lr{ASR}) انتهابه‌انتها گشود \مرجع{gulati2020conformer}. فست کانفورمر (\lr{Fast Conformer}) نیز همین خانواده را با تمرکز بر کاهش هزینهٔ پردازش دنباله‌های بلند توسعه داد \مرجع{rekesh2023fast}. اهمیت این آثار برای مرور حاضر، ارائهٔ گزینه‌های معماری و نشان‌دادن نقش کارایی است. بااین‌حال، مقایسهٔ آن‌ها با مدل‌های ازپیش‌آموخته تنها زمانی معتبر است که داده، روش رمزگشایی و مجموعهٔ آزمون کنترل شوند. در غیر این صورت، برتری مشاهده‌شده ممکن است ناشی از تفاوت در مقیاس پیش‌آموزش یا حجم داده باشد، نه ساختار رمزگذار.

\قسمت{پیکره‌های وب‌مقیاس و نظارت ضعیف}

\زیرقسمت{گردآوری و پالایش گفتار وب}

دو پیکرهٔ «گفتار مردم» (\lr{People's Speech}) و گیگاسپیچ (\lr{GigaSpeech}) نمونه‌های شاخص تبدیل منابع پراکندهٔ وب به داده‌های \lr{ASR} هستند، اما رویکرد یکسانی به گردآوری ندارند. «گفتار مردم» در مقیاسی بزرگ به جست‌وجوی فایل‌های صوتیِ دارای متن همراه و مجوز مناسب پرداخت و سامانهٔ گردآوری را نیز شرح داد \مرجع{galvez2021peoples}. \lr{GigaSpeech} با گردآوری ۱۰ هزار ساعت صوت رونویسی‌شده از کتاب صوتی، پادکست و محتوای ویدئویی، هم‌ترازی اجباری، قطعه‌بندی و پالایش کیفیت را در خط لوله گنجاند \مرجع{chen2021gigaspeech}. مجموعه‌های توسعه و آزمون آن نیز جداگانه و با بازبینی حرفه‌ای تهیه شدند. این دو پروژه نشان می‌دهند که متن همراه هزینهٔ آغاز گردآوری را کاهش می‌دهد، اما ناسازگاری زیرنویس یا ترجمه با صوت، هم‌ترازی و ثبت منشأ را به بخش‌هایی اساسی از ساخت پیکره تبدیل می‌کند.

مقیاس، بازتولیدپذیری و تنوع در این پیکره‌ها هدف‌هایی مرتبط اما گاه متعارض‌اند. جست‌وجوی گسترده تنوع موضوعی را افزایش می‌دهد، ولی تغییر منابع وب بازتولید دقیق را دشوار می‌کند؛ محدودکردن منابع نیز پیگیری منشأ را آسان‌تر می‌سازد، اما ممکن است یک سبک یا مؤسسه را بیش‌ازحد بازنمایی کند. همچنین، کیفیت پذیرفتنی برای گسترش داده‌های آموزشی با کیفیت لازم برای ارزیابی نهایی یکسان نیست: شبه‌برچسب نویزی ممکن است در حجم بالا سودمند باشد، اما مرجع آزمون نویزی منشأ خطا را مبهم می‌کند. ازاین‌رو، ساعت داده تنها پس از تعریف منبع، حذف موارد نامعتبر و قطعه‌بندی معنای روشنی پیدا می‌کند.

در زبان فارسی، نمونه‌ای هم‌سنگ از یک خط لولهٔ بازتولیدپذیر برای تبدیل صوت طولانی و ناهمگون وب به قطعه‌های قابل ممیزی کمتر دیده می‌شود. انتقال مستقیم معیارهای پالایش انگلیسی نیز کافی نیست، زیرا خط فارسی، صورت‌های نوشتاری گوناگون و گفتار محاوره‌ای می‌توانند رابطهٔ متن و صوت را تغییر دهند. بنابراین، گزارش مرحله‌به‌مرحلهٔ گردآوری، قطعه‌بندی، برچسب‌گذاری، نسخه و منشأ هر نمونه به‌اندازهٔ تعداد ساعت‌های نهایی اهمیت دارد.

\زیرقسمت{قطعه‌بندی صوت بلند}

خط لولهٔ \lr{GigaSpeech} قطعه‌بندی را همراه با هم‌ترازی و کنترل کیفیت به کار گرفت تا فایل‌های صوتی طولانی را به واحدهای جمله‌مانند تبدیل کند \مرجع{chen2021gigaspeech}. این رویکرد زمانی مناسب است که متن همراه و زمان‌بندی تقریبی در دسترس باشند. در صوت بدون متن، مرزهای آکوستیکی معمولاً نقطهٔ شروع‌اند. تفاوت این دو وضعیت اهمیت دارد: هم‌ترازی می‌تواند برای یافتن مرزها از متن استفاده کند، درحالی‌که آشکارسازی فعالیت گفتاری (\lr{VAD}) تنها بر شواهد صوتیِ گفتار و ناگفتار تکیه دارد.

هوانگ و همکاران در قطعه‌بند انتهابه‌انتها (\lr{E2E Segmenter}) همین محدودیت را به‌طور مستقیم آزمودند. آن‌ها قطعه‌بندی مبتنی بر \lr{VAD} را با مدلی مقایسه کردند که تعیین مرز را همراه با بازشناسی و با بهره‌گیری از ویژگی‌های غنی‌تر آکوستیکی و زبانی می‌آموخت. آزمایش روی محتوای طولانی و واقعی، شواهد مثبتی به سود قطعه‌بندی آگاه از متن ارائه کرد \مرجع{huang2022e2esegmenter}. اهمیت این نتیجه فراتر از یک معماری مشخص است: هر مکثی نشانهٔ پایان معنایی نیست و مکث ناشی از تردید در میانهٔ جمله ممکن است به‌اشتباه مرز قطعه در نظر گرفته شود. بااین‌حال، این مطالعه بر سامانه و داده‌های انگلیسی در یک سناریوی برخط متمرکز بود؛ بنابراین، نمی‌توان از نتایج آن دربارهٔ کیفیت همهٔ روش‌های \lr{VAD} یا انواع محتوای فارسی نتیجه‌گیری کرد.

در ادبیات قطعه‌بندی، میان سه معیار موازنه برقرار می‌شود: دقت بازشناسی، تأخیر در پایان قطعه و محدودیت طول. قطعه‌های بسیار بلند بافت را حفظ می‌کنند، اما ممکن است از شرایط طولی داده‌های آموزش مدل فراتر روند. در مقابل، قطعه‌های بسیار کوتاه بافت زبانی را می‌گسلند و هزینهٔ رمزگشایی را افزایش می‌دهند. روش‌های آگاه از زبان می‌توانند مرزهای طبیعی‌تری بسازند، اما خود به خروجی مدل وابسته‌اند و در دامنه‌ای ناآشنا ممکن است خطای بازشناسی و خطای تعیین مرز یکدیگر را تشدید کنند. ازاین‌رو، شواهد موجود بر ضرورت ارزیابی مرزها تأکید دارند، نه بر وجود یک قطعه‌بند جهان‌شمول.

\زیرقسمت{خودآموزی و شبه‌برچسب‌های معلم}

کاهن و همکاران خودآموزی را برای \lr{ASR} انتهابه‌انتها با تکیه بر مدل آکوستیکی و زبانی قوی، تولید شبه‌برچسب و پالایه‌های متناسب با خطاهای رمزگشایی بررسی کردند \مرجع{kahn2020selftraining}. آزمایش‌ها نشان دادند که کیفیت شبه‌برچسب و حذف نمونه‌های نامطمئن مستقیماً بر مدل نهایی اثر می‌گذارد، هرچند پالایش بیش‌ازحد نیز داده‌های سودمند را کنار می‌گذارد. بنابراین، صوت بدون برچسب تنها زمانی به نظارت مفید تبدیل می‌شود که متن معلم به‌اندازهٔ کافی به گفتار وفادار باشد.

اسلیم‌آی‌پی‌ال (\lr{slimIPL}) مسئله را از مسیری دیگر دنبال کرد. لیخوماننکو و همکاران وابستگی به مدل زبانی و جست‌وجوی پرهزینه را کاهش دادند و شبه‌برچسب‌ها را هنگام آموزش با یک حافظهٔ پویا بازتولید کردند \مرجع{likhomanenko2021slimipl}. برخلاف خودآموزی یک‌گذره، مدلِ در حال یادگیری به‌تدریج تولید برچسب‌های تازه را بر عهده می‌گیرد. این سازوکار می‌تواند هم‌زمان با بهبود مدل، برچسب‌ها را اصلاح کند، اما ممکن است خطایی را که معتبر پنداشته شده است نیز پیوسته تقویت کند.

این دو مطالعه دربارهٔ نقش مدل زبانی تصویری مکمل ارائه می‌دهند: مدل زبانی قوی در کار کاهن و همکاران کیفیت رمزگشایی را افزایش داد، درحالی‌که \lr{slimIPL} نشان داد با طراحی پایدار می‌توان بدون مدل زبانی خارجی نیز به عملکردی رقابتی رسید. کیفیت شبه‌برچسب در عمل حاصل تعامل میان معلم، دامنه، روش رمزگشایی و قاعدهٔ به‌روزرسانی است؛ ازاین‌رو، نتایج انگلیسی را نمی‌توان بدون ارزیابی مستقل به نام‌های خاص، گفتار محاوره‌ای یا خط فارسی تعمیم داد.

\قسمت{بازشناسی مقاوم و آموزش چندشرایطی}

\زیرقسمت{از نویز عمومی تا ناهمخوانی کانال}

چارچوب آورورا (\lr{AURORA}) امکان مقایسهٔ کنترل‌شدهٔ آموزش با داده‌های پاک و آموزش چندشرایطی را در نویزهای مشخص فراهم کرد و به تثبیت روش ارزیابی مقاومت کمک رساند \مرجع{pearce2000aurora}. لی و همکاران بعدها روش‌های مقاوم‌سازی را در سطوح سیگنال، ویژگی، مدل و آموزش دسته‌بندی کردند و نشان دادند که با وجود پیشرفت معماری‌ها، آموزش در شرایط متنوع همچنان خط مبنایی اساسی است \مرجع{li2014overview}. مزیت این مسیر پژوهشی، امکان نسبت‌دادن تغییر عملکرد به نوع اختلال است؛ محدودیت آن نیز فاصلهٔ میان نویزهای معیار و ترکیب پیچیدهٔ اختلال‌ها در کانال‌های واقعی است.

اسپک‌آگمنت (\lr{SpecAugment}) نمونه‌ای موفق از داده‌افزایی در فضای ویژگی است که با پوشاندن بخش‌هایی از زمان و فرکانس، تعمیم مدل انتهابه‌انتها را بهبود داد \مرجع{park2019specaugment}. بااین‌حال، این تبدیل نه ساختار یک کدک را بازسازی می‌کند و نه رفتار رمزگشا را هنگام افت بسته. ازاین‌رو، داده‌افزایی عمومی و شبیه‌سازی مخابراتی جایگزین یکدیگر نیستند: اولی مدل را در برابر حذف یا تغییر ویژگی‌ها مقاوم‌تر می‌کند و دومی می‌کوشد توزیع مصنوعات مسیر انتقال را بازنمایی کند. در یک مقایسهٔ منصفانه باید روشن شود که بهبود حاصل از افزایش تنوع عمومی داده‌هاست یا از شباهت بیشتر آن‌ها به دامنهٔ هدف ناشی می‌شود.

\زیرقسمت{تلفن، کدک و افت بسته}

سی‌تیمیت (\lr{CTIMIT}) از نخستین پیکره‌هایی بود که گفتار عبورکرده از کانال سلولی را در کنار نسخهٔ مبدأ قرار داد و امکان اندازه‌گیری افت عملکرد ناشی از انتقال را فراهم کرد \مرجع{brown1995ctimit}. مولر و بورلارد نیز با استفاده از یک مدل تحلیلی تلفن، اثر پاسخ فرکانسی، نویز و کدگذاری را به‌طور هم‌زمان بررسی کردند \مرجع{moller2002analytic}. مطالعهٔ نخست بر مشاهدهٔ یک مسیر انتقال واقعی تکیه داشت، درحالی‌که مطالعهٔ دوم امکان جداسازی و کنترل مؤلفه‌ها را افزایش داد. این دو رویکرد در کنار هم نشان می‌دهند که اعتبار بیرونی و کنترل آزمایشگاهی هر دو ضروری‌اند: ضبط واقعی عوامل گوناگون را با یکدیگر درمی‌آمیزد و شبیه‌سازی ممکن است برخی از آن‌ها را در نظر نگیرد.

در انتقال صدا بر بستر پروتکل اینترنت (\lr{VoIP})، دینگ و گوبران اثر نرخ و الگوی افت بسته را بر کیفیت گفتار بررسی کردند و تفاوت میان افت‌های پراکنده و انفجاری را برجسته ساختند \مرجع{ding2003assessment}. فرناندز-گایگو و تولدانو این مسئله را به \lr{ASR} مراکز تماس نزدیک‌تر کردند. آن‌ها کدک کم‌نرخ، چند الگوی افت و روش‌های پنهان‌سازی افت را شبیه‌سازی کردند و مدل آموزش‌دیده با داده‌افزایی را روی داده‌های مصنوعی و ضبط‌های واقعی ارزیابی کردند \مرجع{fernandez2022augmentation}. نقطهٔ قوت مطالعهٔ دوم، پیوند شبیه‌سازی با ارزیابی واقعی است؛ بااین‌حال، یک مجموعه‌دادهٔ مرکز تماس نمی‌تواند نمایندهٔ همهٔ کدک‌ها، شبکه‌ها، زبان‌ها و الگوهای مکالمه باشد.

این ادبیات همچنین هشدار می‌دهد که حذف چند نمونه از شکل موج، معادل افت بسته در هر رمزگشا نیست. پنهان‌سازی افت بسته (\lr{PLC}) ممکن است بسته به حالت پیشین، نوع کدک و طول افت، خروجی متفاوتی تولید کند. تأخیر متغیر، پژواک، هم‌پوشانی گویندگان و قطع‌های طولانی نیز با یک شبیه‌سازی محدود به‌آسانی بازآفرینی نمی‌شوند. بنابراین، نتیجهٔ مثبت روی تخریب مصنوعی تنها دربارهٔ همان توزیع شواهد فراهم می‌کند و باید جدا از نتایج داده‌های واقعی گزارش شود. این تمایز در پژوهش فارسی اهمیت بیشتری دارد، زیرا شواهد منتشرشده دربارهٔ تماس‌های واقعی و معیارهای ارزیابی مشترک اندک‌اند.

برای بازشناسی مقاوم فارسی، این شواهد دو الزام روش‌شناختی ایجاد می‌کنند. نخست، آموزش چندشرایطی باید با یک خط مبنای همسان روی داده‌های عادی مقایسه شود تا اثر تنوع آکوستیکی از تفاوت معماری یا داده جدا بماند. دوم، ارزیابی باید هم تخریب‌های مصنوعیِ کنترل‌شده و هم گفتار واقعی تلفنی یا \lr{VoIP} را در بر گیرد و نتیجهٔ هر شرط را جداگانه گزارش کند. چنین طراحی‌ای امکان می‌دهد انتقال مقاومت از شبیه‌سازی به کانال واقعی بدون پیش‌فرض دربارهٔ نتیجه آزموده شود.

\قسمت{بهسازی گفتار و همجوشی چندنما}

\زیرقسمت{بهسازی سیگنال و دقت بازشناسی}

تفریق طیفی و برآورد کمینهٔ میانگین مربعات دامنهٔ طیفی، پایه‌های روش‌های کلاسیک کاهش نویز را بنا نهادند \مرجع{boll1979spectral,ephraim1984speech}. شبکه‌های عصبیِ پس از آن، از خودرمزگذار عمیق تا حافظهٔ بلند کوتاه‌مدت (\lr{LSTM})، نگاشت پیچیده‌تری میان گفتار آلوده و پاک آموختند \مرجع{lu2013speech,weninger2015speech}. تفاوت اصلی روش‌های جدید با رویکردهای کلاسیک تنها در ظرفیت مدل نیست؛ هدف بهینه‌سازی نیز از مدل‌سازی آماری نویز به بازسازی داده‌محور انتقال یافته است. بااین‌حال، در هر دو خانواده، بهبود معیارهای سیگنالی لزوماً به معنای حفظ نشانه‌های لازم برای بازشناسی نیست.

مور و همکاران بهسازی را هم با معیارهای ابزاری و هم از منظر \lr{ASR} ارزیابی کردند و نشان دادند که نتیجه به سازگاری میان بهساز و بازشناس وابسته است \مرجع{moore2017speech}. ایواموتو و همکاران خطای بهسازی را به نویز باقی‌مانده و مصنوعات پردازشی تفکیک کردند و نشان دادند که مصنوعات ممکن است برای \lr{ASR} مخرب‌تر باشند \مرجع{iwamoto2022bad}. مطالعهٔ بعدی آن‌ها نیز نشان داد که آموزش انتهابه‌انتها با هدف بازشناسی می‌تواند مصنوعات را کاهش دهد، هرچند لزوماً نویز باقی‌مانده را کمتر نمی‌کند \مرجع{iwamoto2024artifacts}. جمع‌بندی این نتایج روشن می‌سازد که بهسازِ مورد استفاده در \lr{ASR} باید بر پایهٔ پیامد زبانی آن سنجیده شود، نه صرفاً بر اساس شباهت شکل موج یا کیفیت ادراکی.

\زیرقسمت{آموزش مشترک و حفظ نمای اصلی}

ونینگر و همکاران یک شبکهٔ \lr{LSTM} بهساز را هم برای بهبود گفتار و هم برای \lr{ASR} مقاوم بررسی کردند و شواهد اولیه‌ای دربارهٔ پیوند هدف بهسازی با وظیفهٔ پایین‌دستی ارائه دادند \مرجع{weninger2015speech}. ژو و همکاران بعدها بهسازی، بازنمایی خودنظارتی و بازشناسی را در چارچوبی مشترک قرار دادند تا اعوجاج ناشی از پیش‌پردازش را با توجه به هدف \lr{ASR} کنترل کنند \مرجع{zhu2022joint}. میزان یکپارچگی در این مطالعات متفاوت است، اما هر دو از جداسازی کامل بهساز و بازشناس فاصله می‌گیرند.

مسیر دیگر، حفظ ورودی اصلی در کنار ورودی بهسازی‌شده است. فوجیموتو و کاوایی ویژگی‌های نویزی و بهسازی‌شده را در یک گذر و با سازوکاری دروازه‌ای ترکیب کردند \مرجع{fujimoto2019onepass}. لی و همکاران، افزون بر دو نمای اصلی و بهسازی‌شده، برآورد نویز را نیز در اختیار مدل آکوستیکی قرار دادند \مرجع{lee2022speech}. در کار ژو و همکاران نیز سازوکار توجه دوگانه، تعامل غنی‌تری میان ویژگی‌های نویزی و بهسازی‌شده ایجاد کرد \مرجع{zhu2022joint}. روند مشترک این مطالعات، حرکت از جایگزینی کامل ورودی به سوی حفظ اطلاعات مکمل است؛ هرچند این رویکرد هزینهٔ معماری و احتمال بی‌استفاده‌ماندن یکی از جریان‌ها را نیز افزایش می‌دهد.

این مطالعات غالباً روی نویز محیطی و با معماری‌های پایهٔ متفاوت انجام شده‌اند. شواهد آن‌ها دربارهٔ مکمل‌بودن نماها را نمی‌توان بی‌واسطه به مصنوعات کدک و افت بسته تعمیم داد، زیرا حذف اطلاعات در یک کانال پارامتری با افزودن نویز ایستا یکسان نیست. افزون بر این، اگر خط مبنای \lr{ASR} از پیش با همان تخریب‌ها و به‌صورت چندشرایطی آموزش دیده باشد، حاشیهٔ بهبود برای بهساز یا همجوشی کوچک‌تر و آزمودن ادعا دشوارتر می‌شود. نبود چنین خط مبنای قدرتمندی می‌تواند سود معماری پیچیده را بیش از اندازه برآورد کند.

تحلیل خودِ سازوکار همجوشی نیز در ادبیات یکدست نیست. بهبود \lr{WER} نشان نمی‌دهد که دروازه واقعاً انتخابی معنادار میان دو نما انجام داده یا یکی از جریان‌ها را تقریباً کنار گذاشته است. حذف آزمایشی هر نما، گزارش رفتار وزن‌ها در شرایط گوناگون و مقایسه با اتصال ساده می‌تواند ادعای استفاده از اطلاعات مکمل را بیازماید. در پژوهش‌های فارسی و تخریب‌های مخابراتی، این نوع تحلیل هنوز به‌اندازهٔ گزارش امتیاز نهایی مورد توجه قرار نگرفته است.

بر پایهٔ این شواهد، حفظ هم‌زمان نمای اصلی و نمای بهسازی‌شده فرضیه‌ای قابل آزمون برای گفتار مخابراتی فارسی است: نمای بهسازی‌شده ممکن است برخی اختلال‌ها را کاهش دهد و نمای اصلی از حذف نشانه‌های بازشناختی جلوگیری کند. داوری دربارهٔ این فرضیه به مقایسه با خط مبنای چندشرایطیِ قوی، حذف آزمایشی هر نما و تحلیل رفتار سازوکار همجوشی نیاز دارد؛ صرف افزایش پیچیدگی معماری یا بهبود معیار سیگنالی برای تأیید آن کافی نیست.

\قسمت{مرور مقایسه‌ای مطالعات منتخب}

جدول~\ref{tab:literature-comparison} مطالعاتی را که بیشترین ارتباط را با زنجیرهٔ استدلال این پژوهش دارند، از نظر دامنهٔ داده، شیوهٔ ارزیابی و محدودیت تعمیم مقایسه می‌کند. مقصود جدول رتبه‌بندی روش‌ها نیست، زیرا داده‌ها، معماری‌ها و معیارهای ارزیابی یکسان نیستند؛ جدول نشان می‌دهد هر دسته از آثار کدام بخش از مسئله را پشتیبانی می‌کند و چه بخشی را برای گفتار مخابراتی فارسی بی‌پاسخ می‌گذارد.

\begin{table}[htbp]
\centering
\caption{مقایسهٔ مطالعات منتخب مرتبط با پژوهش حاضر}
\label{tab:literature-comparison}
\begingroup
\footnotesize
\setlength{\tabcolsep}{3pt}
\renewcommand{\arraystretch}{1.25}
\begin{tabular}{|p{0.12\textwidth}|p{0.19\textwidth}|p{0.18\textwidth}|p{0.17\textwidth}|p{0.22\textwidth}|}
\hline
\textbf{محور} & \textbf{مطالعات نمونه} & \textbf{داده یا دامنه} & \textbf{شیوه یا معیار ارزیابی} & \textbf{محدودیت برای پژوهش حاضر} \\
\hline
فارسی و انتقال & فارس‌دَت، نویسا و سازگارسازی \lr{Whisper} \مرجع{bijankhan1994farsdat,sameti2008nevisa,rahimipour2024persian} & گفتار فارسیِ عمدتاً خوانده‌شده یا پیکره‌محور & دقت یا خطای بازشناسی پس از آموزش و سازگارسازی & پوشش محدود گفتار واقعی تلفنی، کدک و افت بسته \\
\hline
پیکرهٔ وب & «گفتار مردم» و \lr{GigaSpeech} \مرجع{galvez2021peoples,chen2021gigaspeech} & گفتار بلند انگلیسی از منابع وب و چند دامنه & ساعت قابل‌استفاده، کیفیت هم‌ترازی و مجموعهٔ آزمون بازبینی‌شده & تفاوت زبان و خط؛ دشواری بازتولید و ردیابی منابع پویا \\
\hline
شبه‌برچسب & خودآموزی و \lr{slimIPL} \مرجع{kahn2020selftraining,likhomanenko2021slimipl} & صوت بدون برچسب، عمدتاً انگلیسی & \lr{WER} پس از تولید، پالایش یا بازتولید شبه‌برچسب & وابستگی به کیفیت معلم، دامنه و قاعدهٔ به‌روزرسانی \\
\hline
مقاومت مخابراتی & \lr{AURORA}، \lr{CTIMIT} و داده‌افزایی مراکز تماس \مرجع{pearce2000aurora,brown1995ctimit,fernandez2022augmentation} & نویز مصنوعی، کانال سلولی و تماس واقعی & خطای بازشناسی در نویز، کدک و الگوهای افت بسته & محدودیت زبان و کانال؛ ضرورت جداسازی آزمون مصنوعی و واقعی \\
\hline
بهسازی و همجوشی & حفظ نماهای نویزی و بهسازی‌شده و آموزش مشترک \مرجع{weninger2015speech,fujimoto2019onepass,zhu2022joint,iwamoto2024artifacts} & گفتار نویزی و بهسازی‌شده، عمدتاً با نویز محیطی & \lr{WER}، معیارهای سیگنالی و تحلیل مصنوعات بهسازی & تفاوت با تخریب مخابراتی و نیاز به خط مبنای چندشرایطیِ قوی \\
\hline
\end{tabular}
\endgroup
\end{table}

\قسمت{شکاف‌های پژوهشی}

مرور پیشینه چند شکاف مرتبط را آشکار می‌کند. نخست، انتقال چندزبانه برای فارسی امیدبخش است، اما هنوز معلوم نیست آموزش چندشرایطی با شبیه‌سازی پهنای‌باند، کدک و افت بسته تا چه اندازه به گفتار واقعی تلفنی و \lr{VoIP} تعمیم می‌یابد. این پرسش به مقایسه‌ای همسان میان خط مبنای عادی و چندشرایطی و نیز گزارش جداگانهٔ شرایط مصنوعی و واقعی نیاز دارد.

دوم، معلوم نیست حفظ و همجوشی نمای اصلی و بهسازی‌شده، فراتر از یک خط مبنای چندشرایطیِ قوی، برای تخریب‌های مخابراتی فارسی ارزش افزوده ایجاد کند. آزمون این فرضیه باید افزون بر امتیاز نهایی، حذف آزمایشی نماها و رفتار وزن‌های همجوشی را نیز در بر گیرد تا استفادهٔ واقعی از اطلاعات مکمل روشن شود.

سوم، پروژه‌های وب‌مقیاس انگلیسی امکان ساخت پیکره‌های بزرگ را نشان داده‌اند، اما برای فارسی هنوز به زنجیره‌ای بازتولیدپذیر برای گردآوری، اعتبارسنجی منبع، قطعه‌بندی صوت طولانی، نسخه‌بندی برچسب و ثبت فراداده نیاز است. تعداد ساعت‌ها به‌تنهایی کیفیت پیوند صوت--متن یا امکان بازسازی مجموعه را نشان نمی‌دهد.

در نهایت، این شکاف‌ها به یک نیاز مشترک می‌رسند: ارزیابی سامانه‌های فارسی باید هم داده‌محور و هم کنترل‌شده باشد. کیفیت و استقلال داده‌ها، خط مبنای مقاوم و تحلیل ارزش افزودهٔ معماری، سه لایهٔ جدا اما وابسته‌اند. اگر هر لایه جدا از دو لایهٔ دیگر گزارش شود، نسبت‌دادن بهبود یا افت عملکرد به عامل واقعی دشوار خواهد بود. این جمع‌بندی مسیر پژوهش را مشخص می‌کند، بی‌آنکه نتیجهٔ یک روش یا معماری را از پیش مفروض بگیرد.

\قسمت{جمع‌بندی}

پژوهش‌های پیشین نشان می‌دهند که پیش‌آموزش چندزبانه می‌تواند کمبود داده‌های برچسب‌خورده را کاهش دهد، اما نیاز به سازگارسازی و ارزیابی در دامنهٔ هدف را برطرف نمی‌کند. پیکره‌های وب‌مقیاس نیز اهمیت خط لولهٔ قابل ردیابی، قطعه‌بندی و پالایش را روشن کرده‌اند و خودآموزی، در کنار ظرفیت افزایش داده، خطر گسترش خطای معلم را به همراه دارد. از دید مدل‌سازی، آموزش چندشرایطی خط مبنایی جدی برای تخریب‌های مخابراتی است و بهسازی یا همجوشی تنها در مقایسه با چنین مبنایی قابل ارزیابی‌اند. شکاف اصلی برای زبان فارسی، پیوند این ملاحظات داده‌ای و مدلی در ارزیابی‌ای بازتولیدپذیر و بدون پیش‌داوری دربارهٔ نتیجه است.
